\documentclass[11pt]{article}
\usepackage[margin=1in]{geometry}
\usepackage{hyperref}
\usepackage{listings}
\usepackage{xcolor}
\usepackage{booktabs}
\usepackage{enumitem}
\usepackage{parskip}

\lstset{
  basicstyle=\ttfamily\small,
  backgroundcolor=\color{gray!10},
  frame=single,
  breaklines=true,
  columns=fullflexible,
}

\title{\textbf{TCSS 487 Cryptography}\\[0.4em]
       \large Practical Project --- Cryptographic Library \& App --- Part 1}
\author{Thienan Nguyen}
\date{May 2026}

\begin{document}
\maketitle

% ─────────────────────────────────────────────────────────
\section{Overview}

This report describes Part~1 of the cryptographic library and application
project for TCSS~487. The project is written entirely in Java with no external
cryptographic libraries. All symmetric services are built on the SHA-3/SHAKE
(Keccak) sponge construction as specified in FIPS~202.

The submission contains three source files:

\begin{itemize}[noitemsep]
  \item \texttt{SHA3SHAKE.java} --- implementation of the SHA-3/SHAKE sponge.
  \item \texttt{Main.java} --- command-line application offering the four
        required services.
  \item \texttt{App.java} --- empty placeholder class (no logic).
\end{itemize}

% ─────────────────────────────────────────────────────────
\section{Attribution}

The Keccak-f[1600] permutation (round-constant table \texttt{RC},
rho-rotation offsets \texttt{ROTC}, and pi lane-permutation table
\texttt{PILN}) in \texttt{SHA3SHAKE.java} is ported, with minor adaptations,
from Markku-Juhani O.\ Saarinen's public-domain C reference implementation
\emph{tiny\_sha3}:
\begin{center}
  \url{https://github.com/mjosaarinen/tiny_sha3/blob/master/sha3.c}
\end{center}
The sponge wrapper, byte/lane plumbing, padding logic, and the Java public
interface are original work.

% ─────────────────────────────────────────────────────────
\section{Compiling the Project}

All source files are located in the \texttt{src/} directory. Open a terminal,
navigate to \texttt{src/}, and run:

\begin{lstlisting}
javac SHA3SHAKE.java Main.java
\end{lstlisting}

The program is then invoked as:

\begin{lstlisting}
java Main <command> [arguments...]
\end{lstlisting}

where \texttt{<command>} is one of \texttt{hash}, \texttt{mac},
\texttt{encrypt}, or \texttt{decrypt}.

% ─────────────────────────────────────────────────────────
\section{Services}

\subsection{Service 1 --- Cryptographic Hash (SHA-3-256 and SHA-3-512)}

\textbf{Command syntax:}
\begin{lstlisting}
java Main hash <file>
\end{lstlisting}

\begin{itemize}[noitemsep]
  \item \texttt{<file>} --- path to the file to be hashed.
\end{itemize}

\textbf{What it does:}
Reads the entire file into a byte array and computes both the SHA-3-256 and
SHA-3-512 hashes using the static \texttt{SHA3SHAKE.SHA3()} helper. Both
digests are printed to standard output in hexadecimal.

\textbf{Example:}
\begin{lstlisting}
java Main hash test.txt
SHA-3-256: 91468dc1dec55577d293c47a4c092dea71a43c49a5e79064a3e0e055acedc41d
SHA-3-512: 1f6341857b0c8ae0e59e8aff3a2cad25be1f3cd6b339b260425f8bf2585b2876...
\end{lstlisting}

% ─────────────────────────────────────────────────────────
\subsection{Service 2 --- Authentication Tag / MAC (SHAKE-128 and SHAKE-256)}

\textbf{Command syntax:}
\begin{lstlisting}
java Main mac <file> <passphrase> <tagLengthBytes>
\end{lstlisting}

\begin{itemize}[noitemsep]
  \item \texttt{<file>} --- path to the file to authenticate.
  \item \texttt{<passphrase>} --- the secret passphrase (UTF-8 string).
  \item \texttt{<tagLengthBytes>} --- desired MAC output length in bytes
        (e.g.\ 32 for a 256-bit tag).
\end{itemize}

\textbf{What it does:}
Initializes a \texttt{SHA3SHAKE} sponge at each security level (128 and 256),
absorbs the passphrase followed by the file contents, and squeezes the
requested number of bytes. Both tags are printed in hexadecimal.

\textbf{Example:}
\begin{lstlisting}
java Main mac test.txt "mysecret" 32
SHAKE-128 MAC (32 bytes): a3ef5eb8a0e417c9b21953fb1ef1fe392f85039bb9945e65d701897610cc329a
SHAKE-256 MAC (32 bytes): b26b6ee780da70653d85ee28366c160101844f7b631c9a6aca7051292d86d2b8
\end{lstlisting}

% ─────────────────────────────────────────────────────────
\subsection{Service 3 --- Symmetric Encryption}

\textbf{Command syntax:}
\begin{lstlisting}
java Main encrypt <inputFile> <outputFile> <passphrase>
\end{lstlisting}

\begin{itemize}[noitemsep]
  \item \texttt{<inputFile>} --- path to the plaintext file.
  \item \texttt{<outputFile>} --- path where the cryptogram will be written.
  \item \texttt{<passphrase>} --- the secret passphrase (UTF-8 string).
\end{itemize}

\textbf{What it does:}
\begin{enumerate}[noitemsep]
  \item Derives a 128-bit symmetric key:
        \texttt{key = SHAKE-128(passphrase, 128~bits)}.
  \item Generates a random 128-bit (16-byte) nonce using
        \texttt{java.security.SecureRandom}.
  \item Initializes a SHAKE-128 sponge, absorbs the nonce then the key, and
        squeezes a keystream equal in length to the plaintext.
  \item XORs the keystream with the plaintext to produce the ciphertext.
  \item Writes the cryptogram as: \textbf{nonce (16 bytes) $\|$ ciphertext}.
\end{enumerate}

\textbf{Example:}
\begin{lstlisting}
java Main encrypt test.txt test.enc "mysecret"
Encrypted: test.txt -> test.enc
\end{lstlisting}

% ─────────────────────────────────────────────────────────
\subsection{Service 4 --- Symmetric Decryption}

\textbf{Command syntax:}
\begin{lstlisting}
java Main decrypt <inputFile> <outputFile> <passphrase>
\end{lstlisting}

\begin{itemize}[noitemsep]
  \item \texttt{<inputFile>} --- path to the cryptogram produced by
        \texttt{encrypt}.
  \item \texttt{<outputFile>} --- path where the recovered plaintext will be
        written.
  \item \texttt{<passphrase>} --- the same passphrase used during encryption.
\end{itemize}

\textbf{What it does:}
\begin{enumerate}[noitemsep]
  \item Re-derives the 128-bit symmetric key:
        \texttt{key = SHAKE-128(passphrase, 128~bits)}.
  \item Reads the first 16 bytes of the cryptogram as the nonce; the remainder
        is the ciphertext.
  \item Initializes a SHAKE-128 sponge, absorbs the nonce then the key, and
        squeezes a keystream equal in length to the ciphertext.
  \item XORs the keystream with the ciphertext to recover the plaintext.
  \item Writes the recovered plaintext to \texttt{<outputFile>}.
\end{enumerate}

\textbf{Example:}
\begin{lstlisting}
java Main decrypt test.enc test.dec "mysecret"
Decrypted: test.enc -> test.dec
\end{lstlisting}

% ─────────────────────────────────────────────────────────
\section{Implementation Notes}

\subsection*{SHA3SHAKE class}

The class models a plain, non-duplexing sponge over the Keccak-f[1600]
permutation. The rate and capacity are determined by the \texttt{suffix}
argument to \texttt{init()}:

\begin{center}
\begin{tabular}{lrr}
\toprule
\textbf{Mode} & \textbf{Rate (bytes)} & \textbf{Capacity (bytes)} \\
\midrule
SHAKE-128                  & 168 &  32 \\
SHA-3-224                  & 144 &  56 \\
SHA-3-256 / SHAKE-256      & 136 &  64 \\
SHA-3-384                  & 104 &  96 \\
SHA-3-512                  &  72 & 128 \\
\bottomrule
\end{tabular}
\end{center}

SHA-3 finalization uses domain byte \texttt{0x06}; SHAKE finalization uses
\texttt{0x1F}. The pad10*1 rule is applied at the end of the absorb phase
before the first permutation of the squeeze phase.

\subsection*{Keystream cipher}

The stream cipher used for encryption and decryption is SHAKE-128 seeded with
the nonce followed by the derived key. Because SHAKE is an extendable-output
function, it produces an arbitrarily long pseudorandom stream suitable for
use as a stream cipher.

% ─────────────────────────────────────────────────────────
\section{Known Bugs and Limitations}

\begin{itemize}[noitemsep]
  \item The \texttt{mac} command accepts the tag length in \emph{bytes}.
        Passing a non-integer value will cause a \texttt{NumberFormatException}.
  \item Passphrases containing spaces must be quoted on the command line
        (e.g.\ \texttt{"my secret"}).
  \item Decrypting a file not produced by this program's \texttt{encrypt}
        command will produce garbled output without warning, since no MAC
        verification is included in the non-bonus version.
  \item Very large files are read entirely into memory; files larger than the
        available heap will cause an \texttt{OutOfMemoryError}.
\end{itemize}

% ─────────────────────────────────────────────────────────
\section{References}

\begin{itemize}[noitemsep]
  \item National Institute of Standards and Technology.
        \textit{SHA-3 Standard: Permutation-Based Hash and Extendable-Output
        Functions}. FIPS PUB 202, August 2015.
        \url{https://doi.org/10.6028/NIST.FIPS.202}
  \item M.-J.\ O.\ Saarinen. \textit{tiny\_sha3} (public-domain C reference
        implementation). \url{https://github.com/mjosaarinen/tiny_sha3/blob/master/sha3.c}
\end{itemize}

\end{document}
